\newcommand{\derT}[3][d]{\frac{\mathrm{#1}#2}{\mathrm{#1}#3}} % Derivata totale
\newcommand{\derTiL}[3][d]{\mathrm{#1}#2/\mathrm{#1}#3}       % Derivata totale in linea
\newcommand{\derP}[2]{\frac{\partial#1}{\partial#2}}	      % Derivata parziale

\DeclareMathOperator\VolOp{Vol}
\newcommand{\Vol}[1]{\VolOp\left(#1\right)}

\newcommand\de{\mathrm{d}}
\DeclareMathOperator\diag{diag}
\DeclareMathOperator\trace{tr}

\newcommand{\SpazMate}[3]{% Spaz[iatura] Mate[matica]
    \thinmuskip=#1mu\relax
    \medmuskip=#2mu\relax
    \thickmuskip=#3mu\relax
}
\SpazMate{1}{2}{3} % Seleceting math spacing

\newcommand\DeltaD{\Delta^{\mathbf D}}

\newcommand\RPlus{\mathbb R^+}
\newcommand\RPlusStar{\mathbb R_*^+}
\newcommand\NPlus{\mathbb N^+}


\newcommand\FKdiscrete{\teqref{eqFKDiscrete}{[FK]$_d$}}
\newcommand\HDdiscrete{\teqref{eqHDDiscrete}{[HD]$_d$}}
\newcommand\SMOdiscrete{\teqref{eqSMODiscrete}{[SMO]$_d$}}

\newcommand\FKcontinuous{\teqref{eqFKContinuous}{[FK]$_c$}}
\newcommand\HDcontinuous{\teqref{eqHDContinuous}{[HD]$_c$}}
\newcommand\SMOcontinuous{\teqref{eqSMOContinuous}{[SMO]$_c$}}

\newcommand{\TotM}{M_{\text{tot}}}
\newcommand{\TotP}{P_{\text{tot}}}
\newcommand\bOne{\boldsymbol1}

\newcommand\mInf{m_\infty}
\newcommand\MInf{M_\infty}
\newcommand\PInf{P_\infty}

\newcommand\Mass{\text{M}}
% \newcommand\Year{\text{a}}
\DeclareSIUnit\year{a}
% \newcommand\Molarity{\text{M}}
\DeclareSIUnit\molarity{M}

\newcommand\selezionaLingua[2]{%
  \iflanguage{italian}{#1}% Se la lingua corrente è l'inglese usa #1
  {% Altrimenti se la lingua è in italiano usa #2
    \iflanguage{british}{#2}%
    {%
      Linguaggio non considerato % Se non è né italiano né inglese restituisce inglese
    }%
  }%
}

\newcommand\didascalia[1]{%
    \captionsetup{aboveskip=-.2cm}
    \addcontentsline{lof}{figure}{\protect\numberline{\thefigure}#1}%
    \caption*{\figurename~\thefigure: #1}
}

\newcommand\sottoDidascalia[1]{%
    \refstepcounter{subfigure}%
    (\thesubfigure) #1%
}

\newcommand\dimTesto[1]{%
    \pgfmathparse{#1*1.2}%
    \fontsize{#1 pt}{\pgfmathresult pt}\selectfont%
}

\makeatletter
\newcommand{\includetikzgroupplotfigure}
    {\@ifstar{\includetikzgroupplotfigure@star}
    {\includetikzgroupplotfigure@nostar}}

\newcommand{\includetikzgroupplotfigure@nostar}[5]{%
    \begin{figure}[#1]%
        \centering%
        \refstepcounter{figure}%
        \resizebox{#2}{!}{#3}%
        \didascalia{#4}%
        \label{#5}%
    \end{figure}%
}

\newcommand{\includetikzgroupplotfigure@star}[6]{%
    \begin{figure*}[#1]%
        \vspace{#2}%
        \centering%
        \refstepcounter{figure}%
        \resizebox{#3}{!}{#4}%
        \didascalia{#5}%
        \label{#6}%
    \end{figure*}%
}
\makeatother

\newcommand\ie{\textit{i.e.}}

\newcommand\Brain{\mathcal B}
\newcommand\STCyl{\mathcal Q} % Space-time cylinder

\newcommand\vColla[2]{\vspace{#1 plus #2 minus #2}}
\newcommand\hColla[2]{\hspace{#1 plus #2 minus #2}}

\newcommand\HadProd{\smash{\vcenter{\hbox{\scalebox{0.7}{$\odot$}}}}}

\newcommand\HeII{\mathrm{He}_2}